\begin{center}
    {\large \textbf{Giorno 12} \textbar\ Scimmie! Scimmie! Scimmie! \textbar\ 20 Agosto 2024 \textbar\ \textbf{Ubud}, Bali}
\end{center}

\noindent\rule{\textwidth}{0.4pt}
\begin{figure}[h!] % Portrait images below
    \centering
    % First portrait image (on the left)
    \begin{minipage}[h]{0.48\textwidth} % Set width equal to half the page
        \centering
        \includegraphics[width=\textwidth]{images/flags/Screenshot Ubud.png}
    \end{minipage}
    \hfill
    % Text
    \begin{minipage}[h]{0.48\textwidth}
        \raggedright
        \textbf{Superficie}: 4.567 km² \\
        \vspace{0.3cm}
        \textbf{Popolazione}: 3,3 milioni \\
        \vspace{0.3cm}
        \textit{L'isola}
    \end{minipage}
\end{figure}

\landscapeLargeMargin{images/day12/PXL_20240820_001013269.jpg}

\landscapeLargeMargin{images/day12/PXL_20240820_002353735.jpg}

\landscapeLargeMargin{images/day12/PXL_20240820_002936097.jpg}

\landscapeLargeMargin{images/day12/PXL_20240820_003601999.jpg}

\landscapeLargeMargin{images/day12/PXL_20240820_022253276.MP.jpg}

\landscapeLargeMargin{images/day12/PXL_20240820_023041529.jpg}

\landscapeLargeMargin{images/day12/PXL_20240820_023118704.MP.jpg}

\landscapeLargeMargin{images/day12/PXL_20240820_024328604.jpg}

\landscapeLargeMargin{images/day12/PXL_20240820_025921083.PORTRAIT.jpg}

\landscapeLargeMargin{images/day12/PXL_20240820_030517690.PORTRAIT.jpg}

\landscapeLargeMargin{images/day12/PXL_20240820_031409031.jpg}

\landscapeLargeMargin{images/day12/PXL_20240820_033759575.PORTRAIT.jpg}

\noindent 20 Agosto 2024 \newline

Non ci sono moschee nei paraggi del nostro hotel, quindi la notte procede con un sonno tranquillo e continuo. La mattina siamo svegliati dal vociare allegro dei bambini e scopriamo così di essere vicini a una scuola elementare.

La sera precedente abbiamo chiesto di consumare la colazione nella zona comune e quando usciamo dalla stanza ci troviamo davanti alla cameriera, con le nostre bevande calde su un vassoio, puntualissima. La seguiamo fino alla zona oltre la piscina e scegliamo di sederci a dei tavolini di marmo rotondi. Siamo gli unici ospiti, fatta eccezione per un'altra coppia silenziosa.

La nostra colazione comprende una bevanda calda: caffè nero di Bali per me, ottima scelta, e cioccolata calda di soia per Elisa, pessima scelta. Una bevanda fredda: frullato misto di frutta per me, ottima scelta, e Kombucha per Elisa — bevanda altamente frizzante ottenuta dalla fermentazione del tè zuccherato — pessima scelta. Un piatto principale: Chai spiced caramelized banana crepes, ossia crepes alla banana per me, ottima scelta, e Pumpkin French toast \& vanilla syrup per Elisa, che sarebbe stata un'ottima scelta non fosse che consisteva in una singola fetta di toast tagliata a metà con un cucchiaio di marmellata — pessima scelta. Ma tanti elogi a Elisa per aver sperimentato: io dopo la giornata di ieri proprio non me la sentivo e sono andato sul sicuro.

A proposito della giornata di ieri: trascorro la colazione con lo sguardo perso nel vuoto ad ascoltare la diagnosi del mio malessere in barca, frutto di ricerche online che Elisa ha svolto durante la notte. Sostiene che quello che ho avuto sia stato un colpo di calore e in effetti i sintomi tornano tutti. Madre perdoname por mi vida loca, ma potrebbe avere ragione? La crepe alla banana perde un po' di sapore a causa di tutto ciò, ma oggi è il nostro primo giorno a Ubud e per fortuna mi distraggo in fretta.

Ci godiamo la colazione e, fatta eccezione per la Kombucha, è tutto squisito. Dopo aver fatto il pieno di caffè dedico qualche minuto a esplorare la libreria che avevo scoperto la sera prima. A giudicare dalla scelta dei libri, gli ospiti medi di questo resort sembrano essere americani — moltissimi libri in inglese di David Baldacci — o asiatici e nord-europei, visti i titoli in lingua.

Mentre Elisa finisce di prepararsi esco in avanscoperta alla ricerca di una libreria e di una copia in lingua indonesiana di Cent'anni di solitudine. Trovo una libreria, ma il proprietario mi spiega che per i libri in indonesiano devo cambiare città. Non venitemi poi a dire che questo posto non è commerciale.

Per ovviare alla sconfitta mi aggiro per le vie più nascoste dietro Monkey Street alla ricerca di un autista per il tour. La sera prima abbiamo stilato un piano per concentrare più templi possibili in una sola giornata e adesso devo trovare qualcuno che ci accompagni. Dopo aver provato senza successo da un paio di chioschi, chiedo in quello esattamente di fronte all'entrata del Gayatri. Mentre tutti hanno proposto prezzi altissimi o sostenuto che un tour del genere non fosse fattibile in un giorno solo, qui trovo Ketut — letto K-tú — che sostiene che il mio itinerario sia realizzabile e parte da un prezzo di base decisamente più sensato, sul quale riesco anche a contrattare un po'.

Mentre rivediamo insieme le varie tappe, entro in modalità Antiscam: apro una nota sul telefono e inizio a scrivere davanti a lui tutto quello che ci diciamo. Il discorso, in italiano per semplicità, è stato più o meno questo: "Quindi i templi che vedremo sono questi, confermi? L'ora di partenza è questa, confermi? L'ora di ritorno è indicativamente questa, confermi? Il prezzo per due persone è questo, confermi? Non ci sono extra, neanche in caso di ritardi, confermi?" Gli faccio rileggere la nota e dopo aver ottenuto il suo benestare su tutto gli dico che gli avrei fatto sapere. Sempre bene farsi desiderare un po'.

Quando rientro Elisa è pronta e ci incamminiamo verso la prima tappa della mattinata: la Foresta delle Scimmie. Lungo la strada Elisa scatta foto a qualsiasi cosa veda e in particolare quando troviamo una scimmia che girovaga sui tetti le facciamo un vero e proprio servizio fotografico. Ho già il sentore che presto di scimmie non ne potremo più.

Arriviamo alla foresta e facciamo velocemente i biglietti: essendo ancora presto i turisti sono pochi.

\begin{center}
    % Decorative line
    \noindent\rule{0.6\textwidth}{0.4pt}
\end{center}

\textit{Un po' di storia. La Foresta Sacra delle Scimmie si trova nel villaggio di Padang Tegal, Ubud. Fu fondata a partire dall'XI secolo, quando Bali era sotto l'impero del re Sri Aji Jaya Pangus. Il re e la sua famiglia inizialmente utilizzarono l'area come terreno di caccia. Col passare del tempo le comunità locali vicine iniziarono a trasformare parte del terreno in aree residenziali e agricole e in risposta il re emanò un decreto speciale per preservare la foresta e tutti i suoi abitanti selvatici. Gli abitanti dei villaggi circostanti la considerano ancora oggi un importante centro spirituale, economico ed educativo.}

\begin{center}
    % Decorative line
    \noindent\rule{0.6\textwidth}{0.4pt}
\end{center}

Appena varcata la soglia vediamo che le scimmie sono tantissime e che interagiscono con le persone: soprattutto i cuccioli sono abituati a salire sulle spalle dei turisti e ad arrampicarsi. Siamo tentati anche noi di fare una foto con una scimmia sulle spalle, ma i cartelli di pericolo di contagio di rabbia e vaiolo ci dissuadono.

Percorriamo un sentiero nella foresta e ci imbattiamo presto nella grotta di Durgama. Secondo la credenza indù balinese, dopo la morte l'anima continua il suo viaggio verso questa grotta, il cui nome deriva da Durga Mata-Ji, termine per rispettare la dea Durga. All'interno si trova la statua Lingga, simbolo della dea Siwa nella sua forma astratta, e sculture della dea Durga nella sua rappresentazione come figura tantrica con un'espressione arrabbiata. All'ingresso della grotta si trovano invece stupende incisioni murali tratte dalla favola di Tantri Kamandaka: la parete sinistra illustra il processo di raccolta, la processione rituale per accogliere l'arrivo degli ospiti, mentre quella destra illustra il processo di liberazione, il corteo per ricondurli da dove provenivano.

Tutta la foresta sarebbe piena di storia, ma riporto solo i dettagli che ci hanno colpito di più — in particolare quelli relativi all'antico cimitero balinese. Questo cimitero viene usato solo temporaneamente per seppellire le salme fino al momento della cerimonia di cremazione di massa, che nel villaggio di Padang Tegal si tiene ogni cinque anni: in quel giorno tutti i cadaveri sepolti vengono prelevati dalla tomba e cremati secondo la tradizione indù balinese. Affascinante, ma la parola che rende meglio l'idea è "creepy" — soprattutto quando vediamo le masse di terra sollevate in corrispondenza dei luoghi di sepoltura.

Continuiamo a passeggiare nella giungla, incrociando continuamente scimmie. Dappertutto ci sono lavandini con sapone per lavarsi le mani, probabilmente un lascito del COVID. Le scimmie vengono spesso attirate in luoghi specifici dalle guardie e dagli inservienti del parco che distribuiscono loro enormi tuberi, banane, uova e altre cibarie conservate in gabbie di ferro.

Quando arriviamo nei pressi del Tempio Centrale assistiamo a un inserviente che distribuisce alle scimmie uova sode, banane, sacchetti pieni di riso e persino una noce di cocco. Quest'ultima viene presa da una delle scimmie maschio più grandi che, pazientemente, usando denti e zampe sbuccia completamente la parte esterna fino a lasciare solo la noce del frutto. Una volta terminato questo lungo procedimento — ci impiega circa mezz'ora, radunando un vasto pubblico di turisti — spacca la noce sul pavimento di pietra, ne beve l'acqua e infine l'abbandona, così che anche le scimmie più piccole possano cibarsi. Tutta la scena è ipnotica.

\begin{center}
    % Decorative line
    \noindent\rule{0.6\textwidth}{0.4pt}
\end{center}

\textit{Il Tempio Centrale, in cui non ci è permesso entrare, fu fondato nell'anno 1350. È il luogo di culto del dio Shiva, personificazione del "Sang Hyang Widhi Wasa": l'ordine divino, l'unità divina, il dio supremo dell'induismo indonesiano. In passato veniva utilizzato anche per adorare la dea Durga come divinità principale della setta Bhairawa, che un tempo fiorì a Bali ma ora è stata abbandonata, lasciando un'influenza ancora presente nella vita religiosa balinese.}

\begin{center}
    % Decorative line
    \noindent\rule{0.6\textwidth}{0.4pt}
\end{center}

Bella la storia dei sassi, ma siamo qui per le scimmie quindi raccontiamo anche un fatto su di loro. 

Nella Foresta delle Scimmie vicino in totale 1260 scimmie, divise in dieci gruppi, ognuno dei quali possiede un territorio di origine: gruppo Temple, gruppo Selatan, gruppo New Forest, gruppo As Ram, gruppo Centrale, gruppo Est, gruppo Michelin, gruppo Utara, gruppo Cimitero e gruppo Atap.

Ormai stiamo girando da qualche ora e con la colazione leggera del mattino abbiamo fame, ma tirare fuori del cibo qui equivale a essere assaliti, quindi ci rifugiamo in bagno per ingurgitare le provviste. Io assaggio un onigiri di riso e tonno comprato il giorno prima ma ne avanzo metà: terribile. La merenda non ci soddisfa così ci spostiamo in un caffè all'interno della foresta, il cui ingresso è cinto da gabbie di ferro con serrature per impedire alle scimmie di entrare e saccheggiare. Qui assaggiamo un croissant con le mandorle e un dolce a base di frolla con spuma di cocco e crema — forse il dolce più buono che abbia assaggiato in Indonesia. Recuperate le energie usciamo a terminare il giro.

Vediamo le gabbie dove sono state rinchiuse, per protezione, le scimmie che non possono più vivere in natura poiché menomate. Ci raccogliamo in un minuto di silenzio per Bapang, Lanan e Nelson, divenute cieche a causa del morso di un serpente, e per Tumsist, che ha perso una mano e un piede.

Ormai stanchi ci dirigiamo verso l'uscita, ma siamo interrotti su un ponte di legno da una lotta per il territorio tra due gang di scimmie che si saltano addosso e si ringhiano ferocemente. Uno spettacolo che ci intrattiene a lungo e di cui scattiamo decine di foto e video.

La fame non è stata placata dai dolci e appena usciti andiamo alla ricerca di un posto che cucini Ramen, dal momento che abbiamo sentito che qui a Ubud sia molto buono. Troviamo un ristorante giapponese tranquillo, lontano da Monkey Street, dove assaggiare Gyoza e Ramen preparati come da tradizione. Non avremmo dato due lire al posto ma le recensioni erano positive e possiamo confermare l'alta qualità delle materie prime e del servizio, dimostrato da due cameriere perennemente in piedi in fondo al locale pronte a scattare non appena alzavamo lo sguardo.

Terminato il pranzo — e preso un gelato Mochi per Elisa da portare via — torniamo velocemente in hotel per un rapido rifornimento d'acqua e per confermare il tour dell'indomani a Ketut. Durante il tragitto assistiamo a una scena abbastanza surreale: approfittando di una porta socchiusa, una scimmia entra in un negozio di alimentari, ribalta degli stand di dolci ed esce con in mano due banane. I due commessi la sgridano e chiudono la porta, mentre questa — comodamente seduta su un motorino — sbuccia e divora entrambe le banane alla velocità della luce.

Arrivati di fronte al Gayatri non trovo Ketut, ma mi faccio dare il suo numero da un ragazzo che dice di conoscerlo. Scoprirò solo l'indomani che era proprio Ketut! Gli confermo l'intenzione di fare il tour ma decido di sostituire il famoso tempio di Lempuyang con quello dei pipistrelli di Goa Lawah. Ancora non sappiamo cosa ci aspetta. Gli mando la nota preparata al mattino perché tutto rimanga scritto e con quest'ultima modifica riesco a trattare ancora un po' sul prezzo.

Poco male aver rimosso Lempuyang: questo tempio è famoso per le foto su Instagram realizzate con uno specchio posto alla base della fotocamera per simulare un lago che riflette le persone, e i turisti impiegano fino a tre ore di coda solo per essere immortalati in quella buffonata.

\begin{center}
    % Decorative line
    \noindent\rule{0.6\textwidth}{0.4pt}
\end{center}

Si è fatto primo pomeriggio e chiamiamo un Gojek per andare a visitare il tempio di Goa Gajah. Ammetto che Gojek, come era stato per Uber in Messico, funziona molto bene: posso chattare e mandare foto all'autista per indicargli dove venirci a prendere. Ubud e in particolare Monkey Street sono affollatissime, ma appena usciamo dalla città il viaggio procede tranquillo e in una ventina di minuti siamo al tempio. Originariamente volevamo andarci a piedi — circa un'ora di camminata — ma abbiamo desistito per il caldo e per non stremarci ulteriormente.

Di fronte al tempio i venditori ambulanti sono insistenti ma ancora sostenibili. In tanti vogliono venderci il Sarong, il tipico drappo quadrato indossato a mo' di gonna, che a quanto pare è obbligatorio per accedere al tempio. Obbligatorio lo è davvero, ma sappiamo che li prestano gratuitamente all'entrata — quindi anche oggi niente scam. 

Acquistati i biglietti scendiamo le scale verso una vasta area all'aperto. 

\begin{center}
    % Decorative line
    \noindent\rule{0.6\textwidth}{0.4pt}
\end{center}

\textit{Goa Gajah è un tempio pubblico costruito nell'XI secolo, importante centro di meditazione per i sacerdoti provenienti da altri templi. L'ingresso alla grotta è costituito da un enorme elefante intagliato nella roccia, la cui bocca ne è il portale. All'interno si trova una scultura di Ganesha, dio della scienza con aspetto di elefante, un enorme albero piantato durante la costruzione del tempio e tre piscine dove uomini, donne e bambini si purificano prima della preghiera.}

\begin{center}
    % Decorative line
    \noindent\rule{0.6\textwidth}{0.4pt}
\end{center}

L'ingresso della grotta è molto fotogenico e per fortuna all'interno della fontana non assistiamo alle scene caotiche di turisti indisciplinati che ci saremmo aspettati. A parte questo il tempio nell'insieme non è niente di che. Passeggiamo per l'area limitrofa dove si trovano diversi gazebo di preghiera ricchi di offerte lasciate agli dei, poi il nostro giro termina brevemente. Si poteva fare di più.

Usciti dall'area turistica ci incamminiamo per la strada principale del paesino di Bedulu alla ricerca di Samuan Tiga, altro tempio indù. I marciapiedi sono come sempre dissestati e i negozietti sono baracche aperte sulla strada che vendono prodotti di dubbia provenienza e qualità, ma almeno stiamo vivendo la Bali locale. Niente negozi Ralph Lauren su queste stradine.

Arrivati a Samuan Tiga siamo accolti dal custode intento a spazzare per terra: un simpatico volontario piccolo e con gli occhiali che ci chiede una donazione per entrare e ci presta i Sarong da indossare, particolarmente interessato a sapere se ne abbiamo già acquistati di personali. Questo secondo tempio è magnifico, nulla a che vedere col precedente.

\begin{center}
    % Decorative line
    \noindent\rule{0.6\textwidth}{0.4pt}
\end{center}

\textit{Il tempio Samuan Tiga fu costruito durante il regno del re Chandrasangka Warmadewa nel X secolo ed era il tempio reale dell'antica dinastia Warmadewa. Il nome è significativo: "samuan" significa "incontro" e "tiga" significa "tre". A differenza dei tre cortili classici di altri templi balinesi, Samuan Tiga è costituito da sette padiglioni collegati da scale che conducono al cortile più interno — la sala conferenze degli dei e dei santi. Il tempio è noto anche per due rituali rari: il Siyat Sampian, in cui decine di devote donne in stato quasi incosciente si attaccano a vicenda usando composizioni di foglie di cocco, seguito da non meno di trecento uomini; e il Sanghyang Jaran, una danza in trance in cui il sacerdote del tempio, posseduto da una divinità ancestrale, assume il movimento di un cavallo e danza tra i carboni ardenti.}

\begin{center}
    % Decorative line
    \noindent\rule{0.6\textwidth}{0.4pt}
\end{center}

Ovviamente non assistiamo a nessun rito, ma possiamo esplorare i vari livelli del tempio in totale tranquillità: non c'è praticamente nessuno, solo un paio di visitatori sul punto di andarsene. Riempiamo di ricordi gli occhi e i rullini (le schede SD) e infine ci avviamo verso l'uscita, dove ritroviamo il custode.

Attacchiamo bottone con lui: voglio sapere come mai ci sia così poca gente. Mi racconta che da inizio giornata ha avuto solo sei visitatori, che lui lì è un volontario e che purtroppo i fondi scarseggiano. Parla a fatica inglese e lo capisco a stento, ma tra le cose che riesco a cogliere c'è che "ogni business ha il suo tempio." Scopriamo anche il motivo del suo grande interesse verso i nostri Sarong: sua moglie li vende e dopo la nostra visita andrà ad accompagnarla al mercato. È molto gentile e ci chiede se ne vogliamo acquistare uno da lui, ma Eli ne ha già decine a casa regalati dai suoi genitori e io so che finirebbe in un cassetto a vita. Con dispiacere gli diciamo di no e lo salutiamo. Usciti dal tempio ci sediamo di fronte a un negozietto ad aspettare il Gojek e nell'attesa ci ripassa davanti in motorino, salutandoci sempre col sorriso.

Ubud come al solito è congestionata e il viaggio di ritorno dura quasi il doppio dell'andata. Si è fatta ora di cena e dopo aver posato gli zaini al Gayatri usciamo alla ricerca di un ristorantino. Almeno quella è l'idea, ma Elisa è "rapita dalla magia del luogo" — almeno così lo definisce lei, mentre mi redarguisce mentre scrivo questa pagina di diario: "tu cercavi di spegnere il mio entusiasmo con la tua obiettività e razionalità." Sarà, ma per me è stato come parlare al telefono con qualcuno che non sente una parola di quello che dici e ti parla sopra in continuazione.

Prima di cena riusciamo a comprare anelli e braccialetti da una bancarella, dove ho sparato un prezzo bassissimo come base per contrattare e mi sono sentito rispondere subito di sì — segno che ho sbagliato le stime in pieno. E poi tutta una serie di cucchiai, cucchiaini e cucchiaioni di conchiglie e legno che odio con tutto l'odio che uno può dedicare a delle posate. Usciti finalmente dal negozio di posate — vendeva anche altro, ma per me resterà sempre un negozio di posate — siamo andati nel posto più vicino che facesse un hamburger. È stata una sorta di rivincita personale.

Il Tropical Restaurant Ubud era davvero carino: completamente aperto, con musica dal vivo, decorazioni in bambù e la tipica atmosfera dei locali della zona. La sistemazione però non è stata delle migliori, visto che ci hanno fatto accomodare a un bancone con vista sul vicolo. Poco male: il mio hamburger è andato giù alla perfezione insieme a un frullato di mela e menta, mentre Elisa ha optato per il suo immancabile Mie Goreng vegetariano accompagnato da una birra Bintang.

Finita la cena non ci siamo trattenuti a lungo: il giorno dopo ci attendeva un tour prenotato per le sei del mattino e dovevamo ancora organizzarci per la colazione. Per fortuna il Gayatri ci è venuto in soccorso, ricordandoci con un messaggio di prenotarla. Ho colto subito l'occasione chiedendo una colazione da asporto e ci hanno proposto delle breakfast box da ritirare al mattino prima di partire — un altro punto a loro favore.

Per il pranzo serviva un'altra soluzione. Avremmo potuto fermarci lungo il tour, rischiando però di tornare troppo tardi, così abbiamo deciso di anticipare. Usciti dal locale abbiamo attraversato per l'ennesima volta Monkey Street e siamo tornati al Coco Market, il supermercato visitato la sera prima. Qui abbiamo preso il necessario per i panini: pane, pomodori e mozzarella. Detto così sembra un ottimo menù, ma preciso che la mozzarella era sottovuoto e a prima vista sembrava quasi una toma.

Con la spesa fatta siamo tornati in hotel, dove Elisa si è messa a preparare i panini mentre io, ormai sfinito, sono crollato a letto.


\pagestyle{empty} % disable numbers

\landscapeLargeMargin{images/day12/PXL_20240820_003451668.jpg}

%coppia
\splitImageLeft{images/day12/PXL_20240820_003931457.jpg}
\splitImageRight{images/day12/PXL_20240820_003931457.jpg}
%

%coppia
\landscapeLargeMargin{images/day12/PXL_20240820_022551690.jpg}
\landscapeLargeMargin{images/day12/PXL_20240820_022640683.PORTRAIT.jpg}
%

%%coppia
\fourPortraitImages{images/day12/PXL_20240820_022845408.PORTRAIT.jpg}
                   {images/day12/PXL_20240820_024752015.PORTRAIT.jpg}
                   {images/day12/PXL_20240820_030329802.jpg}
                   {images/day12/PXL_20240820_063211386.jpg}
\twoImagesLargeMargins{images/day12/PXL_20240820_033714984.jpg}
                      {images/day12/PXL_20240820_025921083.PORTRAIT.jpg}
%

%coppia
\fourPortraitImages{images/day12/PXL_20240820_031409031.jpg}
                   {images/day12/PXL_20240820_031606936.jpg}
                   {images/day12/PXL_20240820_032528000.jpg}
                   {images/day12/PXL_20240820_032538463.jpg}
\fourPortraitImages{images/day12/PXL_20240820_041822571.jpg}
                   {images/day12/PXL_20240820_042242469.jpg}
                   {images/day12/PXL_20240820_043529452.jpg}
                   {images/day12/PXL_20240820_043829397.jpg}
%

%coppia
%%% TODO: images/day12/PXL_20240820_050510940.MP.jpg
%  Qui foto di Eli
%

%coppia
\splitImageLeft{images/day12/PXL_20240820_080511036.jpg}
\splitImageRight{images/day12/PXL_20240820_080511036.jpg}
%

%coppia
\threeImagesThinMargins{images/day12/PXL_20240820_080419633.MP.jpg}
                       {images/day12/PXL_20240820_081329954.jpg}
                       {images/day12/PXL_20240820_081354466.MP.jpg}
\twoImagesLargeMargins{images/day12/PXL_20240820_082143602.MP.jpg}
                      {images/day12/PXL_20240820_082214989.MP.jpg}
%

%coppia
\twoImagesLargeMargins{images/day12/PXL_20240820_075812063.MP.jpg}
                      {images/day12/PXL_20240820_081046764.jpg}
\portraitFullscreen{images/day12/PXL_20240820_085242215.jpg}
%

%coppia
\landscapeThinMargin{images/day12/PXL_20240820_090858069.jpg}
\portraitFullscreen{images/day12/PXL_20240820_091532100.jpg}
%

%coppia
\landscapeThinMargin{images/day12/PXL_20240820_091744221.PORTRAIT.jpg}
\twoImagesLargeMargins{images/day12/PXL_20240820_085946969.jpg}
                      {images/day12/PXL_20240820_090109391.jpg}
%

%coppia
\twoImagesLargeMargins{images/day12/PXL_20240820_090154081.jpg}
                      {images/day12/PXL_20240820_090227226.jpg}
\landscapeThinMargin{images/day12/PXL_20240820_090419407.jpg}

%coppia
\splitImageLeft{images/day12/PXL_20240820_093559375.MP.jpg}
\splitImageRight{images/day12/PXL_20240820_093559375.MP.jpg}
%

\landscapeLargeMargin{images/day12/PXL_20240820_114739192.jpg}

\pagestyle{fancy} % enable numbers